%% Six full proofs for the separately specified general finite-CNF CDCL core.
\begin{annexitem}{annex:p12-cdcl-uip-local-sound}{thm:p12-cdcl-uip-local-sound}{Reason and resolution invariants}
\begin{proof}
A consistent trail assigns each variable at most once.  If a clause is unit,
its sole unassigned literal must be true in every model extending the trail,
because all its other literals are false there.  Appending that literal
therefore preserves every such model.  There are finitely many variables and
each propagation assigns a new one, so repeated unit propagation terminates.
A falsified clause prohibits every extending model.  An empty clause does
so immediately; a tautological clause cannot be falsified or unit under a
consistent assignment.  The empty database imposes no constraint.

Every recorded reason was unit before its propagated literal was appended.
All variables of its other literals were already assigned, so its incoming
graph edges strictly follow trail order.  After cancellation of a suffix of
levels, every remaining reason uses only retained earlier entries.  Keeping
the reason clause itself is necessary to use this fact in later analysis.

For resolution, let the parents be \(A\cup\{x\}\) and
\(B\cup\{\neg x\}\), where the designated pivot occurrences are removed
from \(A,B\).  Any valuation satisfying both parents satisfies \(A\cup B\):
if it makes \(x\) false, the first parent supplies a true literal of \(A\);
if it makes \(x\) true, the second supplies a true literal of \(B\).
If the resolvent is tautological it is already true under every valuation;
if it is empty these two cases show the parents have no common model.
Applying this argument to any model of the input proves that resolving two
input consequences produces another input consequence.  Induction on
learning steps, and inside each step on its resolution sequence, establishes
entailment of every learned clause.  Adding such a clause removes no input
model, and every enlarged-database model still satisfies the retained input.
Thus the model set is exactly preserved.  Unit propagation has the narrower
trail-extension meaning above; decisions are not themselves entailed facts.
\end{proof}
\end{annexitem}

\begin{annexitem}{annex:p12-cdcl-uip-analysis}{thm:p12-cdcl-uip-analysis}{First-UIP analysis, assertion, and backjumping}
\begin{proof}
At a positive-level conflict, at least one conflict-clause variable has the
current level.  Otherwise the clause was already false in the earlier retained
prefix, contradicting the required completed propagation before the current
decision.  The same conclusion applies after a backjumped learned assertion:
the inserted clause is unit and then true, and any ensuing new conflict must
depend on a newly asserted or propagated current-level assignment.

Initialize \(\omega_L\) to that falsified clause.  Suppose it has more
than one current-level literal and choose its latest true complement \(p\)
on the trail.  It cannot be the current decision, because that is the first
entry of the level and another current-level entry would be later.  It is
therefore propagated and has the retained reason \(\alpha(p)\).  That
reason contains \(p\); all its other literals were false when \(p\)
was appended and remain false in the conflict trail.  Resolving removes
\(\neg p\) from the temporary clause and \(p\) from the reason.
All remaining literals are false in the same consistent trail.  They cannot
contain both polarities of a variable.  Entailment follows from
Annex~\ref{annex:p12-cdcl-uip-local-sound}.

Each newly introduced current-level variable precedes \(p\), and every
other current-level variable already in \(\omega_L\) also precedes it.
Thus the greatest current-level trail position strictly decreases.  A reason
of a propagated current-level assignment above level zero has a current-level
predecessor.  For an old clause, if all false antecedents had lower levels,
it would already have been unit at the preceding propagation boundary.  For
the newly inserted asserting clause, its maximal nonasserting level equals
the enqueueing level by construction.  Subsequent propagations inherit the
same property.  Consequently elimination never removes the last current-
level frontier entry.  The finite decreasing position measure proves that
the procedure stops with exactly one such entry, rather than getting stuck
at zero entries or requiring the reason of a decision.

Here is the domination argument, including why the stopping node is the
\emph{first} UIP.  Interpret each temporary false literal as the node of its
true complement.  Initially these frontier nodes are exactly the predecessors
of the conflict.  Replacing a propagated frontier node by its reason
predecessors preserves a cut of all paths from the current decision to the
conflict.  Current-level graph nodes are reachable from that decision, by
the current-predecessor property just proved.  Edges cannot run from the
current level to an older level, since levels are nondecreasing along the
block-ordered trail.  Therefore, when only a single current-level frontier
node remains, every decision-to-conflict path passes through it.  This node
is a UIP.

For each current-level frontier node there is also a path from it to the
conflict whose internal current-level nodes have already been eliminated.
This is initially immediate and is preserved when a reason predecessor is
inserted: prepend its edge to the eliminated node's witness path.  The
latest-position pivot rule eliminates nodes in strictly decreasing trail
order.  If a UIP \(z\) lies downstream of the eventual stopping node,
consider the earlier stage at which \(z\) is the latest current-level
frontier node.  Such a stage must occur: reverse expansion cannot bypass
\(z\) on every conflict path, and the eventual node is upstream of it.
All previously eliminated current nodes at that stage are later than \(z\).
Any other current frontier node would be earlier than \(z\), and its witness
path to the conflict would use only those already eliminated later nodes.
A path from the current decision to this earlier node cannot contain the
later \(z\), by acyclicity and trail order.  Combining it with the witness
path would give a decision-to-conflict path avoiding \(z\), contradicting
its domination.  Hence that earlier stage had only \(z\) in its current
frontier, and the algorithm would already have stopped there.  There is no
downstream UIP beyond the actual stopping node, proving the first-UIP claim.

Let \(a\) be the final clause's unique current-level literal.  Every other
literal is false and has level at most \(b<d\).  Cancellation to \(b\)
removes the variable of \(a\) and retains all others false.  Hence the
clause is unit and enqueuing \(a\) is justified.  If \(b>0\), cancelling
to any lower level also removes at least one other literal of level \(b\),
so at least two literals are unassigned and the clause is no longer unit.
If \(b=0\), no lower permitted level exists.  Thus this is the lowest
retained assertion level.  A learned unit has no other literals and asserts
at zero.  Root literals in a larger learned clause remain false at zero and
need not be deleted from it.  A conflict at zero is handled by root
refutation, not by a nonexistent positive-level decision or reason.
\end{proof}
\end{annexitem}

\begin{annexitem}{annex:p12-cdcl-uip-watch-correct}{thm:p12-cdcl-uip-watch-correct}{Watched propagation and level-prefix backtracking}
\begin{proof}
We prove the event invariant needed for both forward search and rollback.
Every nonunit live clause has two distinct watches and exact list membership.
Here a pending event includes a queued event or a still-unvisited clause
reference in the watch list currently being processed; dequeueing its literal
does not make those references already processed.
Every false watched literal either has an unprocessed event in the current
level's queue, or has the other watch true at a level no greater than its own.
At a conflict the just-processed conflicting watches are an allowed terminal
exception; both belong to the level that will be erased, unless the conflict
is at zero and search terminates.  Completed propagation has no pending
event, so each nonunit clause either has both watches nonfalse or has a true
watch paired with a false watch at an equal or greater level.

Input initialization satisfies the invariant: before assignments all watches
are unassigned.  Root units have their separate enqueue path; an empty clause
conflicts immediately.  At insertion of a learned nonunit clause, its asserting
watch is made true at level \(b\), while its second watched literal was
chosen false at the same maximal level \(b\).  The required level inequality
therefore holds.  A learned unit is separately retained and enqueued at zero.
These facts include insertion into a nonempty trail rather than assuming all
attachments happen at an empty assignment.

Consider processing a watched false literal \(l\).  Its assignment was made
in the current propagation phase, so its level is the current level; no new
decision can have occurred while its event was pending.  If its other watch
\(q\) is true, \(\delta(q)\leq\delta(l)\), and keeping the pair
establishes the invariant for this processed event.  If a replacement \(r\)
is nonfalse, moving the watch preserves two distinct positions and the exact
old/new list membership.  If \(q\) is false, its event must still be pending.
Indeed, had it already been processed and left in place, it would have had a
true companion whose level was no greater than its own; that companion
cannot have become false during a consistent forward phase.  A prior
processed event with no true companion would instead have moved that watch,
enqueued its companion true, or reported a conflict.  The remaining false
\(q\) is consequently protected by its pending event.  If the replacement
is true, it also has level no greater than the current \(q\) event.  If it
is unassigned, it is not false and imposes no new event obligation.

When no replacement exists, every nonwatched literal is false.  If \(q\)
is unassigned, the clause is genuinely unit, so enqueueing it is sound.  Its
new level equals that of the false event, providing the needed true companion.
If \(q\) is false, all literals of the clause are false and the reported
conflict is genuine.  A binary clause is exactly this no-replacement case.
A tautological clause cannot produce a false or unit condition under a
consistent assignment, so it is harmless if retained; complementarity never
licenses treating its two occurrences as a unit.  Repeated true enqueues
add no events, and opposite enqueues report a real incompatible consequence.
Restoring unprocessed list references on conflict and detaching only deleted
clauses preserves membership for the next phase.

There are at most \(\lvert X\rvert\) new assignments in one propagation
phase, hence that many new queue entries.  Each event scans a finite watch
list and each replacement search scans a finite clause.  The phase terminates.
If it finishes conflict-free, take any nonunit active clause.  With no pending
false watch, either it has a true watched literal and is satisfied, or both
watches are unassigned, since nonfalse watches without a true one must be
unassigned.  The latter supplies two distinct unassigned positions, so that
clause is neither false nor unit.  Separately enqueued unit clauses are true.
Therefore no required unit or conflict has been missed at the fixed point.
Together with soundness of every enqueue this proves propagation correctness;
it does not assert that every clause's complete satisfied/unresolved status
is maintained by the lazy data structure.

Finally cancel whole levels to \(b\).  A nonfalse watch can become only
unassigned or remain at its former value, never newly false.  A false watch
whose assignment is erased becomes unassigned.  A false watch whose assignment
is retained had either a pending current-level event or a true companion at
no greater level.  In ordinary conflict cancellation all pending current-level
assignments are erased because the jump is below that level; a restart is at
a completed-propagation boundary and has none.  Thus a retained false watch
has a true companion at level at most its own and therefore at most \(b\);
the companion is retained as well.  Watch moves during an erased phase do
not change this argument, since every performed move satisfied the same
event/level invariant.  At a nonroot conflict the terminal offending watches
are newly false current-level entries and are erased; an older retained false
watch cannot have lost its older true companion without that companion also
having been erased earlier.  Hence the post-cancellation watches already
satisfy the completed invariant.  Nothing needs restoration.  The new learned
assertion then has the insertion treatment above.  If cancellation were within
a level, or if the learned second watch had lower level than its assertion,
the companion implication would fail: the body's wrong-watch example exhibits
exactly that omitted condition.
\end{proof}
\end{annexitem}

\begin{annexitem}{annex:p12-cdcl-uip-search-correct}{thm:p12-cdcl-uip-search-correct}{CDCL soundness and a terminating complete policy}
\begin{proof}
Every retained learned clause is entailed by the immutable input, by the
resolution invariant.  Since every input clause is retained, the active
database and the input have exactly the same total models.  Deleting only
entailed learned clauses cannot change this equality.  A restart or backjump
changes the candidate trail, not the input or the entailment of surviving
clauses.  Locking the still-used reasons ensures that the premise clauses
needed by the analysis construction remain available.  The watched
propagation theorem ensures the loop's conflicts and assertions are genuine.

At a SAT return, the trail is total and the propagation phase has no conflict.
Every input clause is therefore true: if one were false, it would be an active
false clause prohibited by the watched fixed-point result.  At a root UNSAT
return, induct on root propagations.  Their initial unit reasons and subsequent
reasons are input consequences, and their other literals are false by root
facts already entailed by the input.  Every root fact is consequently entailed
by the input.  The root-conflict clause is also entailed, but all its literals
are false under those facts, so no input model exists.  An initially empty
clause is the zero-step instance.  These prove soundness independently of
how many restarts or permitted deletions precede a return.

For liveness with retained learning, we show that every nonroot conflict
learns a clause not already active.  The final asserting clause is false on
the conflict trail and has exactly one current-level variable.  If it were an
already active unit clause, its level-zero fact would be true, a contradiction.
If it were an already active nonunit clause, consider its two distinct
watches under the event invariant proved in the preceding annex.  A false
watch from an older level has no pending event: propagation was completed
before the present level began.  Its companion must therefore still be true
at a no-greater level, contradicting falsification of the whole clause.
Accordingly both false watches of an already active falsified clause must
be current-level entries.  This also covers the just-reported terminal
conflict exception, whose offending watches are both current-level entries.
Two distinct watch positions then give two distinct current-level variables
in that non-tautological clause, contrary to the learned clause's assertion
property.  Thus the learned clause cannot already be active.  This argument
applies equally to clauses inserted after an earlier decision: their immediate
assertion and maximal-level second watch establish the same event invariant.
It uses actual propagation and attachment correctness, not an assumption
that an asserting clause is new merely because it has been given a new name.

There are at most \(3^{\lvert X\rvert}\) non-tautological literal-set
clauses: each variable is absent, positive, or negative.  Every learned clause
uses these fixed variables and is non-tautological, so retaining every learned
clause allows only finitely many nonroot conflicts.  Between conflicts, a
propagation phase is finite, and each decision assigns a previously unassigned
variable.  With no voluntary restart there can be at most \(\lvert X\rvert\)
decisions without a conflict or a total assignment.  Hence an infinite run
would require infinitely many conflicts, contradicting the finite-clause
bound.  The loop terminates.  Its only returns are the two sound terminal
outcomes, so it is also complete for both satisfiable and unsatisfiable inputs.

If restarts occur only after completed nonroot learning steps, each restart
is charged to one of the finitely many conflicts.  There are finitely many
such restarts and the same finite-phase argument applies.  The stated absence
of an infinite administrative delay rules out stuttering forever without
performing the next required propagation or decision.  Arbitrary voluntary
restarts before progress, or deletion followed by relearning the same clause,
invalidate this counting argument.  The SAT example in the body already
gives an infinite sound run under the former policy.  The proof therefore
does not infer termination from VSIDS, randomization, a solver name, or the
mere existence of a finite decision space.
\end{proof}
\end{annexitem}

\begin{annexitem}{annex:p12-cdcl-uip-activity-scale}{thm:p12-cdcl-uip-activity-scale}{Ideal-real lazy decay preserves priority order}
\begin{proof}
Let \(s_t(x)\) be eager scores after \(t\) completed decays, and
\(a_t(x)\) the stored lazy scores.  Start with equal scores and unit bump
increment.  Maintain \(s_t(x)=\rho^t a_t(x)\) for every variable, with
lazy increment \(\Delta_t=\rho^{-t}\).  For a bump vector \(h_t(x)\),
the eager update is \(s_t(x)+h_t(x)\), and the lazy update is
\(a_t(x)+\Delta_t h_t(x)\); multiplying the latter by \(\rho^t\)
gives exactly the former.  Performing the next eager decay multiplies this
by \(\rho\), while the lazy method leaves scores unchanged and replaces
the increment by \(\Delta_t/\rho\).  Thus the invariant holds at
\(t+1\).  The factor \(\rho^t\) is positive, so comparisons and ties
are identical at all times.  Simultaneously multiplying every stored score
and the increment by any positive constant only changes the common conversion
factor and preserves all future comparisons by the same calculation.  The
proof uses exact reals.  Rounding, overflow and implementation-specific tie
handling are not premises of it and are not claimed verified.
\end{proof}
\end{annexitem}

\begin{annexitem}{annex:p12-cdcl-uip-lookahead-sound}{thm:p12-cdcl-uip-lookahead-sound}{Failed-literal and two-branch consequence soundness}
\begin{proof}
Let \(\nu\) be the retained consistent prefix and let \(v\) be any total
model of the active database extending it.  If the temporary branch \(x\)
has a unit conflict, propagation soundness shows that no such \(v\) can
make \(x\) true.  Therefore every extending model makes \(\neg x\)
true.  If both branches conflict, a Boolean valuation must choose one of
the two polarities, so there is no extending model at all.  If both branches
propagate \(l\), choose the branch matching \(v(x)\).  Soundness of its
propagations shows \(v(l)=1\).  Thus \(l\) is true in every extending
model.  These arguments remain valid when one branch already has no model;
they do not assume propagation finds every semantically impossible branch.

For a global clause, let the true prefix literals be \(p_1,\ldots,p_k\).
The failed branch \(x\) establishes entailment of
\(\neg p_1\lor\cdots\lor\neg p_k\lor\neg x\): a model falsifying
this clause would extend the prefix and make \(x\) true, which was excluded.
The analogous clause with \(l\) follows for the common consequence.  Dropping
prefix literals without a separate discharge argument strengthens the claim
and is not justified.  Restoring the original trail after each temporary
test, while retaining only such justified global clauses or prefix-relative
assignments, prevents temporary branch decisions from being mistaken for
input consequences.
\end{proof}
\end{annexitem}
